% Agent-authored source transcription and faithful translation, verified against PDF 25-28.
% The sample starts at Section 2.4; it excludes the preceding Section 2.3 and later Box 2.1.
% block: chain-maps-heading page:25 type:section
\subsection{{Chain maps}}\label{sec:chain-maps}

% block: chain-map-definition page:25 type:ordinary-paragraph
{A \emph{chain map} from a chain complex $(V,\partial)$ to a second chain complex $(W,\xi)$ is a sequence of linear maps $f_k:V_k\to W_k$ such that the following diagram commutes:}

% block: chain-map-diagram page:25 type:equation
\[
\begin{tikzcd}[column sep=1.2em,row sep=2.2em]
\cdots \arrow[r] & V_{k+1} \arrow[r,"\partial_{k+1}"] \arrow[d,"f_{k+1}"'] & V_k \arrow[r,"\partial_k"] \arrow[d,"f_k"'] & V_{k-1} \arrow[r] \arrow[d,"f_{k-1}"'] & \cdots \\
\cdots \arrow[r] & W_{k+1} \arrow[r,"\xi_{k+1}"'] & W_k \arrow[r,"\xi_k"'] & W_{k-1} \arrow[r] & \cdots
\end{tikzcd}
\]

% block: chain-map-induced-homology page:25 type:ordinary-paragraph
{In other words, a chain map is a graded linear map of degree 0 which commutes with the differentials $f\circ\partial=\xi\circ f$. Since $f(\partial c)=\xi f(c)$ for $c\in V_{k+1}$, $f$ maps the $k$-boundaries of $V$ to $k$-boundaries of $W$, and the same equation applied to $c\in\mathcal Z_k$ shows that $f$ maps $k$-cycles to $k$-cycles. It follows that the chain map $f$ induces a linear map $\bar f$ on homology, i.e., the map $\bar f_k:\mathcal H_k(V)\to\mathcal H_k(W)$ which takes the homology class $\bar v$ of some $v\in V_k$ to the homology class $\overline{f_kv}$ of $f_kv\in W_k$. This map is functorial. That is, if $(X,\zeta)$ is a third chain complex and $g$ is a chain map from $W$ to $X$, then $g\circ f$ is a chain map from $V$ to $X$ and $\overline{g\circ f}=\bar g\circ\bar f$.}

% block: chain-projection-definition page:25 type:ordinary-paragraph
{A case we shall find of great interest later is when the complex $W$ is a subcomplex of $V$ and the maps $f_k$ are projections of $V_k$ onto $W_k$, i.e., $f_kv=v$ for $v\in W_k$. In this case we say that $f$ is a \emph{chain projection}.}

% block: chain-projection-surjectivity page:25 type:proposition
\begin{Proposition}\label{prop:chain-projection-homology-surjection}
{A chain projection from a complex to a subcomplex induces a surjection on homology.}
\end{Proposition}

% block: chain-projection-surjectivity-proof page:25 type:proof
\begin{Proof}
{The subcomplex inclusion maps $W_k\hookrightarrow V_k$ certainly form a chain map, $i$. When $f$ is a projection, $f\circ i$ is the identity map from $W$ to itself, and so $\bar f\circ\bar i=\overline{f\circ i}$ is the identity on homology and is, in particular, surjective. Therefore $\bar f$ must be surjective.}
\end{Proof}

% block: cochain-complexes-heading page:26 type:section
\subsection{{Cochain complexes}}\label{sec:cochain-complexes}

% block: cochain-complex-definition page:26 type:ordinary-paragraph
{A cochain complex is the same thing as a chain complex except with the indexing reversed: cochain differentials have degree $+1$ rather than $-1$. One typically writes the index of a cochain complex as a superscript rather than a subscript and often uses $d$ rather than $\partial$ to denote the differential:}

% block: cochain-complex-sequence page:26 type:equation
\[
\cdots\to V^{k-1}\xrightarrow{d^{k-1}}V^k\xrightarrow{d^k}V^{k+1}\to\cdots.
\]

% block: cohomology-definition page:26 type:ordinary-paragraph
{Of course, all the definitions and results above carry over in an obvious way. For example, the $k$th cohomology space $\mathcal H^k$ is defined as the quotient of the cocycles $\mathcal Z^k$ by the coboundaries $\mathcal B^k$.}

% block: dual-cochain-complex page:26 type:ordinary-paragraph
{Since any linear map $L:V\to W$ between vector spaces induces a dual map $L':W'\to V'$ between the dual spaces, any chain complex gives rise to a dual cochain complex. For example, from the simplicial chain complex associated to a simplicial complex, we obtain the simplicial cochain complex. Since the $k$-dimensional simplices form a basis for $C_k$, a $k$-cochain is an element of $C^k:=C'_k$, i.e., a function $\omega$ which assigns to each $k$-dimensional simplex a real number $\omega(f)$ (extended to chains by linearity). The differential $\partial^k:C^k\to C^{k+1}$ is defined by $\partial^k=\partial'_{k+1}$, i.e.,}

% block: dual-differential-evaluation page:26 type:equation
\[
(\partial^k\omega)g=\omega(\partial_{k+1}g),\qquad \omega\in C^k,\quad g\in C_{k+1}.
\]

% block: incidence-transpose page:26 type:ordinary-paragraph
{This means that the matrix of the operator $\partial^k$ is the transpose of the oriented incidence matrix which represents $\partial_{k+1}$. For example, $\partial^0$ applied to a vertex is the sum of the edges incident on that vertex, each multiplied by $1$ or $-1$ according to whether the edge points into or out of the vertex. Of course, the information contained in the simplicial cochain complex}

% block: simplicial-cochain-complex page:26 type:equation
\begin{equation}\label{eq:simplicial-cochain-complex}
0\to C^0\xrightarrow{\partial^0}C^1\xrightarrow{\partial^1}\cdots\xrightarrow{\partial^{n-1}}C^n\to0
\end{equation}

% block: cohomology-duality page:26 type:ordinary-paragraph
{is essentially the same as that contained in the simplicial chain complex. In particular, it can be shown that the cohomology space $\mathcal H^k$ is naturally isomorphic to $\mathcal H'_k$, the dual of the corresponding homology space, and so its dimension is the $k$th Betti number.}

% block: de-rham-heading page:26 type:section
\subsection{{Example: The de Rham complex}}\label{sec:de-rham-complex}

% block: de-rham-introduction page:26 type:ordinary-paragraph
{We will have a lot more to say about the de Rham complex}

% block: differential-forms-complex page:26 type:equation
\begin{equation}\label{eq:differential-forms-de-rham-complex}
0\to\Lambda^0(\Omega)\xrightarrow{d^0}\Lambda^1(\Omega)\xrightarrow{d^1}\cdots\xrightarrow{d^{n-1}}\Lambda^n(\Omega)\to0
\end{equation}

% block: de-rham-background page:26 type:ordinary-paragraph
{associated to any $n$-dimensional domain or smooth manifold $\Omega$. Here $\Lambda^k(\Omega)$ is the space of (smooth) differential $k$-forms on $\Omega$ and the operator $d$ is the exterior derivative, as introduced by Élie Cartan [\ref{bib:cartan-differential-expressions}]. We shall put off the introduction of the language of differential forms until Chapter~\ref{chap:exterior-calculus}. For now we just describe the case of $\Omega$ a domain in $\mathbb R^3$, for which the de Rham complex can be described in the terminology of vector calculus.}

% block: vector-calculus-complex-intro page:26 type:ordinary-paragraph
{The complex is then}

% block: vector-calculus-de-rham-complex page:26 type:equation
\begin{equation}\label{eq:vector-calculus-de-rham-complex}
0\to C^\infty(\Omega)\xrightarrow{\operatorname{grad}}C^\infty(\Omega;\mathbb R^3)\xrightarrow{\operatorname{curl}}C^\infty(\Omega;\mathbb R^3)\xrightarrow{\operatorname{div}}C^\infty(\Omega)\to0.
\end{equation}

% block: calculus-identities page:27 type:ordinary-paragraph
{The calculus identities $\operatorname{curl}\operatorname{grad}\phi=0$ and $\operatorname{div}\operatorname{curl}u=0$ ensure that \eqref{eq:vector-calculus-de-rham-complex} is indeed a complex. Its cohomology spaces $\mathcal H^k_{\mathrm{dR}}$ comprise the de Rham cohomology of the domain.}

% block: de-rham-map-zero-cochains page:27 type:ordinary-paragraph
{If the domain $\Omega$ is triangulated by some simplicial complex $\mathcal S$, there is an interesting cochain map from the de Rham complex to the simplicial cochain complex called the \emph{de Rham map}. Namely, we associate to a function $\phi\in C^\infty(\Omega)$ the 0-cochain $F^0\phi$ determined by $F^0\phi(x)=\phi(x)$ for each $x\in\Delta_0(\mathcal S)$. As we shall see in Chapter~\ref{chap:exterior-calculus} a differential $k$-form may be integrated on an oriented $k$-dimensional manifold, and this is the simplest case of $k=0$. This is extended by linearity to arbitrary 0-chains, so $F^0\phi(\sum c_ix_i)=\sum c_i\phi(x_i)$. Similarly we associate to a vector field $\psi\in C^\infty(\Omega;\mathbb R^3)$ the 1-cochain whose value on any oriented edge $e\in\Delta_1(\mathcal S)$ is}

% block: de-rham-edge-integral page:27 type:equation
\[
F^1\psi(e)=\int_e\psi\cdot t\,ds_e,\qquad e\in\Delta_1(\mathcal S),
\]

% block: de-rham-trace-forms page:27 type:ordinary-paragraph
{where the $t$ denotes the unit vector tangent to the edge $e$ directed according to the orientation. The vector field $\psi$ may be viewed as a differential 1-form, and it has a trace on $e$ which is a differential 1-form on the one-dimensional manifold $e$ and so may be integrated over it. The line integral above is the expression of that integral in vector calculus terminology. A vector field in $\mathbb R^3$ may alternately be viewed as a differential 2-form, in which case the integral of its trace over a surface is defined and may be realized as}

% block: de-rham-face-integral page:27 type:equation
\[
F^2\psi(f)=\int_f\psi\cdot n\,ds_f,\qquad f\in\Delta_2(\mathcal S),
\]

% block: de-rham-volume-form-intro page:27 type:ordinary-paragraph
{where now $n$ is the unit normal to the face $f$, its direction determined by the orientation of $f$ according to the right-hand rule. Finally, a scalar function $\phi$ may be thought of as a 3-form, leading to}

% block: de-rham-volume-integral page:27 type:equation
\[
F^3\phi(T)=\int_T\phi\,dx
\]

% block: de-rham-connecting-map page:27 type:ordinary-paragraph
{for a positively oriented tetrahedron $T\in\Delta_3(\mathcal S)$. We have thus defined the de Rham map connecting the de Rham complex \eqref{eq:vector-calculus-de-rham-complex} to the simplicial cochain complex \eqref{eq:simplicial-cochain-complex}:}

% block: de-rham-commuting-diagram page:27 type:equation
\begin{equation}\label{eq:de-rham-commuting-map}
\begin{tikzcd}[column sep=1.1em,row sep=2.2em]
0 \arrow[r] & C^\infty(\Omega) \arrow[r,"\operatorname{grad}"] \arrow[d,"F^0"'] & C^\infty(\Omega;\mathbb R^3) \arrow[r,"\operatorname{curl}"] \arrow[d,"F^1"'] & C^\infty(\Omega;\mathbb R^3) \arrow[r,"\operatorname{div}"] \arrow[d,"F^2"'] & C^\infty(\Omega) \arrow[r] \arrow[d,"F^3"'] & 0 \\
0 \arrow[r] & C^0 \arrow[r,"\partial^0"'] & C^1 \arrow[r,"\partial^1"'] & C^2 \arrow[r,"\partial^2"'] & C^3 \arrow[r] & 0
\end{tikzcd}
\end{equation}

% block: de-rham-first-cell-verification page:27 type:ordinary-paragraph
{The de Rham map is a cochain map, which is to say that this diagram commutes. For example, let us verify that $F^1\circ\operatorname{grad}=\partial^0\circ F^0$. We must show that, for an arbitrary smooth function $\phi$, the 1-cochains $F^1\operatorname{grad}\phi$ and $\partial^0F^0\phi$ coincide, i.e., that}

% block: de-rham-first-cell-identity page:27 type:equation
\[
F^1\operatorname{grad}\phi([x_0,x_1])=\partial^0F^0\phi([x_0,x_1])
\]

% block: oriented-edge-unit-tangent pages:27,28 type:ordinary-paragraph
{for any oriented edge $[x_0,x_1]$ in the complex. Letting $t$ denote the unit vector $(x_1-x_0)/|x_1-x_0|$, we have}

% block: first-cell-calculation page:28 type:equation
\[
\begin{aligned}
F^1\operatorname{grad}\phi([x_0,x_1])
&=\int_{[x_0,x_1]}\operatorname{grad}\phi\cdot t
&&\text{{(definition of $F^1$)}}\\
&=\phi(x_1)-\phi(x_0)
&&\text{{(fundamental theorem of calculus)}}\\
&=F^0\phi([x_1]-[x_0])
&&\text{{(definition of $F^0$)}}\\
&=F^0\phi(\partial_1[x_0,x_1])
&&\text{{(definition of $\partial_1$)}}\\
&=\partial^0F^0\phi([x_0,x_1])
&&(\partial^0=\partial'_1).
\end{aligned}
\]

% block: remaining-cells page:28 type:ordinary-paragraph
{The commutativity of the other two cells in the diagram is proven in the same way, using the Stokes theorem and the divergence theorem from vector calculus:}

% block: stokes-divergence-identities page:28 type:equation
\[
\int_f\operatorname{curl}\psi\cdot n=\int_{\partial f}\psi\cdot t,\qquad
\int_T\operatorname{div}\psi=\int_{\partial T}\psi\cdot n.
\]

% block: comparison-of-complexes page:28 type:ordinary-paragraph
{The de Rham complex which constitutes the first row of the diagram \eqref{eq:de-rham-commuting-map} is very different from the simplicial cochain complex which is its second row. The former is composed of infinite dimensional function spaces connected by differential operators, while the latter consists of finite dimensional spaces with operators that are essentially combinatorial (multiplication by matrices whose entries are all $1$, $-1$, or $0$). However, the cohomology of the two complexes is identical. This statement is a form of de Rham's theorem.}

% block: de-rham-isomorphism page:28 type:theorem
\begin{Theorem}{de Rham}\label{thm:de-rham-isomorphism}
{The map on cohomology induced by the de Rham map is an isomorphism from the de Rham cohomology spaces onto the simplicial cohomology spaces.}
\end{Theorem}

% block: singular-and-simplicial-cohomology page:28 type:remark
\begin{Remark}\label{rem:singular-simplicial-de-rham}
{De Rham's theorem as just stated equates de Rham cohomology, which is constructed analytically using differential operators, with the more geometric simplicial cohomology, constructed from chains of simplices. Although simplicial homology is in some sense the simplest and most easily computed homology theory, topologists often prefer singular homology for its greater generality and theoretical flexibility. In particular, simplicial homology is defined only for spaces which can be triangulated (or are homeomorphic images of such spaces), while singular homology, being built from chains of singular simplices (arbitrary continuous images of simplices) is defined for any topological space. An early, quite elementary, result in most expositions of homology is that the two approaches give identical homology (or cohomology) spaces when both apply, that is, for triangulated spaces. De Rham's theorem, which is a deeper result, is then usually stated as giving an isomorphism between de Rham cohomology and singular cohomology on a smooth manifold.}
\end{Remark}

% block: de-rham-later-proof page:28 type:ordinary-paragraph
{De Rham's theorem is a deep result. However, we shall see in Section~\ref{sec:whitney-forms-de-rham} that a version of it can be proven quite easily once we have assembled the tools of FEEC.}
